\documentclass[aspectratio=169]{beamer}

\usepackage{fontspec}
\usepackage{polyglossia}
\usepackage{color}
\usepackage{listings}
\setdefaultlanguage{russian}
\usepackage{csquotes}
\usetheme{metropolis}
\metroset{background=light}
\metroset{sectionpage=none}
\metroset{numbering=fraction}

\definecolor{ssaublue}{RGB}{77,144,205}
\setbeamercolor{frametitle}{fg=white, bg=ssaublue}
\setbeamercolor{title separator}{fg=ssaublue}

\setmainfont[
    Path = ../fonts/elektratextpro/,
    Extension = .otf,
    UprightFont = elektratextpro,
    BoldFont = elektratextpro_bold,
    ItalicFont = elektratextpro_italic,
    BoldItalicFont = elektratextpro_bolditalic
]{elektratextpro}

\setsansfont[
    Path = ../fonts/elektratextpro/,
    Extension = .otf,
    UprightFont = elektratextpro,
    BoldFont = elektratextpro_bold,
    ItalicFont = elektratextpro_italic,
    BoldItalicFont = elektratextpro_bolditalic
]{elektratextpro}

\setmonofont[
    Path = ../fonts/elektratextpro/,
    Extension = .otf,
    UprightFont = elektratextpro,
    BoldFont = elektratextpro_bold,
    ItalicFont = elektratextpro_italic,
    BoldItalicFont = elektratextpro_bolditalic
]{elektratextpro}

\newfontfamily\cyrillicfont[
    Path = ../fonts/elektratextpro/,
    Extension = .otf,
    UprightFont = elektratextpro,
    BoldFont = elektratextpro_bold,
    ItalicFont = elektratextpro_italic,
    BoldItalicFont = elektratextpro_bolditalic
]{elektratextpro}

\newfontfamily\cyrillicfontsf[
    Path = ../fonts/elektratextpro/,
    Extension = .otf,
    UprightFont = elektratextpro,
    BoldFont = elektratextpro_bold,
    ItalicFont = elektratextpro_italic,
    BoldItalicFont = elektratextpro_bolditalic
]{elektratextpro}

\newfontfamily\cyrillicfonttt[
    Path = ../fonts/elektratextpro/,
    Extension = .otf,
    UprightFont = elektratextpro,
    BoldFont = elektratextpro_bold,
    ItalicFont = elektratextpro_italic,
    BoldItalicFont = elektratextpro_bolditalic
]{elektratextpro}

% Настройка листингов для Beamer
\lstset{
    basicstyle=\ttfamily\footnotesize,
    commentstyle=\color{gray},
    columns=fullflexible,
    keepspaces=true,
    breaklines=false,
    extendedchars=true,
    texcl=true,
    showstringspaces=false
}

\title{Управление конфигурацией}
\date{1 октября 2025 г.}

\begin{document}
    \maketitle

    \begin{frame}{Что мы имеем после прошлых лекций}
        \textbf{Наш CI/CD pipeline работает:}
        
        \begin{itemize}
            \item Git push → автоматический деплой через GitHub Actions
            \item CI/CD pipeline выполняет тесты и обновляет код в основном и тестовом окружении
        \end{itemize}

        \textbf{Но есть проблема:}

        Сервер (окружение) настраивается руками.
    \end{frame}

    \begin{frame}{Почему это проблема?}
        \textbf{Проблемы ручной настройки серверов:}
        
        \begin{itemize}
            \item \textbf{Ошибки} — человеческий фактор при настройке
            \item \textbf{Дрейф конфигурации} — серверы постепенно расходятся в настройках
            \item \textbf{Масштабирование} — нужно 10 серверов? 10 раз повторяем
        \end{itemize}
        
        \vspace{0.5cm}
        
        \textbf{Вопрос:} как автоматизировать настройку серверов?
    \end{frame}

    \begin{frame}{Два подхода к автоматизации}
        \begin{columns}[t]
            \begin{column}[t]{0.48\textwidth}
                \textbf{Baked (образы):}
                \begin{itemize}
                    \item Заранее готовим образ с настроенным ПО
                    \item Разворачиваем готовые серверы из образа
                    \item Быстрый запуск
                \end{itemize}
                
                \vspace{0.3cm}
                \textbf{Примеры:} Docker, Packer и др.
            \end{column}
            
            \begin{column}[t]{0.48\textwidth}
                \textbf{Fry (конфигурация):}
                \begin{itemize}
                    \item Берем чистый сервер
                    \item Автоматически настраиваем его скриптами
                    \item Гибкость в настройке
                \end{itemize}
                
                \vspace{0.3cm}
                \textbf{Примеры:} Ansible, Chef, и др.
            \end{column}
        \end{columns}
        
    \end{frame}

    \begin{frame}{Что такое Ansible}
        \textbf{Ansible} — инструмент для автоматизации настройки серверов
        
        \vspace{0.3cm}
        
        \textbf{Основные особенности:}
        \begin{itemize}
            \item \textbf{Agentless} — не нужно ставить агенты на сервера
            \item \textbf{SSH} — использует обычное SSH соединение
            \item \textbf{YAML} — конфигурация в простом формате
            \item \textbf{Идемпотентность} — можно запускать много раз безопасно
        \end{itemize}
        
        \vspace{0.5cm}
        
        \textbf{Принцип работы:} описываем желаемое состояние сервера в файлах, Ansible приводит сервер к этому состоянию
    \end{frame}

    \begin{frame}{Основные концепции Ansible}
        \textbf{Inventory} — список серверов для управления
        \begin{itemize}
            \item IP адреса или доменные имена
            \item Группировка серверов (веб-сервера, БД, и т.д.)
        \end{itemize}
        
        \vspace{0.3cm}
        
        \textbf{Playbook} — сценарий действий в YAML формате
        \begin{itemize}
            \item Последовательность задач (tasks)
            \item Применяется к группам серверов из inventory
        \end{itemize}
        
        \vspace{0.3cm}
        
        \textbf{Modules} — готовые блоки для типовых задач
        \begin{itemize}
            \item apt, yum — установка пакетов
            \item copy, template — работа с файлами
            \item service — управление сервисами
        \end{itemize}
    \end{frame}

    \begin{frame}{Демонстрация: Настраиваем nginx с Ansible}
        \textbf{Задача:} автоматизировать настройку nginx на сервере
        
        \vspace{0.3cm}
        
        \textbf{Что покажем:}
        \begin{itemize}
            \item inventory файл со списком серверов
            \item Пишем playbook для установки и настройки nginx
            \item Добавляем запуск playbook в workflow
        \end{itemize}

        \vspace{0.2cm}
        \textbf{Демо-репозиторий:} \\
        \href{https://github.com/prafdin/devops-demo-website/tree/cm-demo}{\texttt{https://github.com/prafdin/devops-demo-website/tree/cm-demo}}
    \end{frame}

    \begin{frame}{Что запомнить}
        \begin{itemize}
            \item Ручная настройка серверов — источник проблем и дрейфа конфигурации
            \item Два подхода к автоматизации конфигурации: baked (образы) и fry (конфигурация)
            \item Ansible — один из инструментов для конфигурации сервера
        \end{itemize}

        \vspace{0.5cm}

        \begin{center}
            \textbf{Вопросы?}
        \end{center}
    \end{frame}

\end{document}